\scanpage{33}
% archive sheet number: 32
% slanted
\margin{hyp. de \unsure{l'induction|induction}}
$3^\circ)$ L'hom can
\[
\Omega_X|Y\longrightarrow\underline{\mathrm{Hom}}(\unsure{P|\rho}_{Y/X},\Omega_Y)
\]
est injectif [une orientation de $X$ en $y\in Y$ est connue quand on sait comment \inserted{au vois. de $y$} une rive \inserted{de $Y$\deleted{\unsure{ en $y$}}} détermine une orientation de $Y$]

$4^\circ)$ Pour tte $\omega\in\Gamma(X)$, l'hom correspondant
\[
P_{Y/X}\longrightarrow\Omega_Y
\]
% slanted
\margin{hyp. de non trivialité}
est injectif [deux rives différentes définissent des orientations différentes de $Y$]

$5^\circ)$ Pour tout $X\in\mathrm{Ob}\,\unsure{V|\mathcal{V}}$, $\Omega(X)\to X$ est couvrant (il existe localement des orientations \unsure{$g$|$?$} de $X$; \inserted{ou encore : $\forall n\geq 1$, $\Omega(\mathbb{R}^n)$ est non vide})

% slanted
\margin{\deleted{\unsure{hyp.}}principe de \deleted{\unsure{pseudo}}l'identité}
$6^\circ)$ Pour tt $X\in\mathrm{Ob}\,V$, $\Omega(X)$ interprété comme espace étalé sur $X$ est séparé sur $X$, ou encore : une orientation \unsure{de l'un|de l'espace|d'une} $X\in\mathrm{Ob}\,V$ est connue en un pt quand on la connait en des pts arbitrairement voisins. (Il suffit de l'exiger pour $\mathbb{R}^n$, $n\geq 1$, en fait, il sera vrai à posteriori qu'une orientation de $\mathbb{R}^n$, \inserted{ou de tt $X$ connexe,} est connue quand on la connaît en un seul point\ldots)
